\section{NeuS2}
Chắc chắn rồi, dựa trên bài báo bạn cung cấp về NeuS2, tôi đã tổng hợp thông tin tập trung vào phương pháp, ưu nhược điểm và chuẩn bị một section LaTeX cùng file BibTeX tương ứng.

1. Mã LaTeX cho Section trong Luận văn:

Đoạn mã

\section{NeuS2: Học Nhanh Bề Mặt Ngầm Neural để Tái tạo Đa Góc nhìn}

Các phương pháp tái tạo bề mặt 3D dựa trên biểu diễn ngầm neural, như NeuS \cite{Wang2021NeuS}, đã chứng minh khả năng tạo ra kết quả chất lượng rất cao cho cảnh tĩnh \cite{Wang2023NeuS2_source_1}. Tuy nhiên, một hạn chế lớn của NeuS là thời gian huấn luyện cực kỳ dài (khoảng 8 giờ), khiến nó gần như không thể áp dụng cho các cảnh động với hàng nghìn khung hình \cite{Wang2023NeuS2_source_2, Wang2023NeuS2_source_17}. Mặt khác, các phương pháp tăng tốc như Instant-NGP \cite{Muller2022InstantNGP} tuy nhanh nhưng chất lượng hình học trích xuất từ trường mật độ (density field) lại bị nhiễu và thiếu các ràng buộc bề mặt \cite{Wang2023NeuS2_source_19, Wang2023NeuS2_source_94}. Việc kết hợp trực tiếp NeuS và Instant-NGP gặp thách thức lớn trong việc tính toán hiệu quả đạo hàm bậc hai (cần thiết cho ràng buộc Eikonal của NeuS) khi sử dụng các MLP được tối ưu hóa cao bằng CUDA \cite{Wang2023NeuS2_source_23, Wang2023NeuS2_source_24, Wang2023NeuS2_source_95, Wang2023NeuS2_source_96, Wang2023NeuS2_source_97, Wang2023NeuS2_source_98}.

NeuS2 \cite{Wang2023NeuS2} được đề xuất để giải quyết những vấn đề này, nhằm mục tiêu học nhanh các bề mặt ngầm neural chất lượng cao từ ảnh đa góc nhìn cho cả cảnh tĩnh và động.

\subsection{Phương pháp Mô hình NeuS2}

NeuS2 kế thừa ý tưởng biểu diễn bề mặt bằng hàm khoảng cách có dấu (SDF) và cơ chế volume rendering không lệch (unbiased) từ NeuS \cite{Wang2021NeuS}, đồng thời tích hợp các kỹ thuật tăng tốc từ Instant-NGP \cite{Muller2022InstantNGP} và đề xuất các cải tiến mới.

\subsubsection{Kiến trúc cho Cảnh tĩnh}
\begin{itemize}
    \item \textbf{Biểu diễn kết hợp SDF và Hash Encoding:} Thay vì dùng MLP lớn, NeuS2 sử dụng một MLP \textit{nhỏ} để học hàm SDF $f$. Đầu vào của MLP này là sự kết hợp (concatenate) của tọa độ điểm 3D $x$ và vector đặc trưng $h_\Omega(x)$ được truy vấn từ cấu trúc mã hóa băm đa phân giải (Multi-resolution Hash Encoding) với các mục trong bảng băm $\Omega$ có thể học được \cite{Wang2023NeuS2_source_4, Wang2023NeuS2_source_22, Wang2023NeuS2_source_86, Wang2023NeuS2_source_110, Wang2023NeuS2_source_112}. Mạng SDF này dự đoán giá trị SDF $d$ và một vector đặc trưng hình học $g$ \cite{Wang2023NeuS2_source_112}.
    \item \textbf{Mạng Màu:} Một MLP nhỏ khác dự đoán màu sắc $c$, nhận đầu vào là vị trí $x$, pháp tuyến $n = \nabla_x d$, hướng nhìn $v$, giá trị SDF $d$, và đặc trưng hình học $g$ \cite{Wang2023NeuS2_source_114, Wang2023NeuS2_source_115}.
    \item \textbf{Volume Rendering:} Sử dụng công thức volume rendering không lệch của NeuS \cite{Wang2021NeuS} để tính màu pixel từ SDF và màu dự đoán dọc theo tia nhìn \cite{Wang2023NeuS2_source_116, Wang2023NeuS2_source_405, Wang2023NeuS2_source_409}. Quá trình ray marching được tăng tốc bằng occupancy grid tương tự Instant-NGP \cite{Wang2023NeuS2_source_117, Wang2023NeuS2_source_411, Wang2023NeuS2_source_412}.
    \item \textbf{Tính toán Đạo hàm Bậc hai Hiệu quả:} Đây là đóng góp kỹ thuật quan trọng của NeuS2 để xử lý ràng buộc Eikonal $\mathcal{L}_{eikonal} = \sum (\|\nabla_x d\| - 1)^2$ một cách hiệu quả \cite{Wang2023NeuS2_source_119, Wang2023NeuS2_source_120}. Thay vì dùng finite difference (dễ gây mất ổn định \cite{Wang2023NeuS2_source_99, Wang2023NeuS2_source_178}) hoặc autograd của các framework (chậm \cite{Wang2023NeuS2_source_124}), NeuS2 đưa ra công thức giải tích \textit{chính xác} và rút gọn cho đạo hàm bậc hai của hàm SDF $d$ đối với các tham số của mạng SDF và bảng băm, đặc biệt tận dụng tính chất của MLP dùng ReLU ($\frac{\partial^2 y}{\partial x^2} = 0$) \cite{Wang2023NeuS2_source_4, Wang2023NeuS2_source_25, Wang2023NeuS2_source_32, Wang2023NeuS2_source_100, Wang2023NeuS2_source_125, Wang2023NeuS2_source_129, Wang2023NeuS2_source_130}. Điều này cho phép cài đặt hiệu quả bằng CUDA kernel, tăng tốc đáng kể quá trình backpropagation \cite{Wang2023NeuS2_source_4, Wang2023NeuS2_source_132}.
    \item \textbf{Huấn luyện Tiến bộ (Progressive Training):} Để cải thiện sự ổn định và chất lượng khi dùng hash encoding, NeuS2 áp dụng chiến lược huấn luyện tăng dần độ phân giải. Các mức trong bảng băm được kích hoạt tuần tự từ thô đến mịn trong quá trình huấn luyện thông qua một hệ số trọng số phụ thuộc vào tiến trình huấn luyện \cite{Wang2023NeuS2_source_5, Wang2023NeuS2_source_26, Wang2023NeuS2_source_33, Wang2023NeuS2_source_109, Wang2023NeuS2_source_135, Wang2023NeuS2_source_136, Wang2023NeuS2_source_137}.
\end{itemize}

\subsubsection{Kiến trúc cho Cảnh động}
NeuS2 mở rộng sang tái tạo cảnh động bằng hai thành phần chính:
\begin{itemize}
    \item \textbf{Huấn luyện Tăng cường (Incremental Training):} Thay vì huấn luyện mỗi frame độc lập, NeuS2 tận dụng tính liên tục thời gian. Frame đầu tiên được huấn luyện từ đầu, các frame tiếp theo được khởi tạo từ các tham số (bảng băm, trọng số MLP) đã học của frame trước đó và chỉ cần fine-tune trong thời gian ngắn (~20 giây/frame). Điều này giúp tăng tốc đáng kể quá trình học cho toàn bộ chuỗi video \cite{Wang2023NeuS2_source_6, Wang2023NeuS2_source_28, Wang2023NeuS2_source_103, Wang2023NeuS2_source_140, Wang2023NeuS2_source_142, Wang2023NeuS2_source_143}.
    \item \textbf{Dự đoán Biến đổi Toàn cục (Global Transformation Prediction - GTP):} Khi có chuyển động lớn giữa các frame, việc fine-tune trực tiếp có thể gặp vấn đề (mắc kẹt ở cực tiểu địa phương). NeuS2 đề xuất dự đoán một phép biến đổi cứng (xoay $R_i$, tịnh tiến $T_i$) giữa frame hiện tại $i$ và frame trước đó $i-1$. Các điểm 3D trong frame $i$ được ánh xạ về không gian gốc (canonical space - thường là frame 0) thông qua chuỗi biến đổi tích lũy trước khi được đưa vào mạng SDF/màu. GTP giúp ổn định quá trình huấn luyện tăng cường và cho phép mô hình hóa các chuỗi dài với chuyển động lớn trong một vùng không gian giới hạn, tiết kiệm bộ nhớ \cite{Wang2023NeuS2_source_6, Wang2023NeuS2_source_30, Wang2023NeuS2_source_34, Wang2023NeuS2_source_90, Wang2023NeuS2_source_104, Wang2023NeuS2_source_145, Wang2023NeuS2_source_146, Wang2023NeuS2_source_148, Wang2023NeuS2_source_150, Wang2023NeuS2_source_151, Wang2023NeuS2_source_152}. GTP được tối ưu đồng thời trong quá trình fine-tune \cite{Wang2023NeuS2_source_163}.
\end{itemize}

\subsection{Ưu điểm và Nhược điểm}

\subsubsection{Ưu điểm}
\begin{itemize}
    \item \textbf{Tốc độ Vượt trội:} NeuS2 nhanh hơn NeuS khoảng hai bậc độ lớn (ví dụ: 5 phút so với 8 giờ cho cảnh tĩnh trên DTU) \cite{Wang2023NeuS2_source_3, Wang2023NeuS2_source_10, Wang2023NeuS2_source_50, Wang2023NeuS2_source_157}. Tái tạo cảnh động chỉ mất khoảng 20 giây cho mỗi frame (sau frame đầu tiên) \cite{Wang2023NeuS2_source_11, Wang2023NeuS2_source_21, Wang2023NeuS2_source_197, Wang2023NeuS2_source_208}.
    \item \textbf{Chất lượng Tái tạo Cao:} Đạt được chất lượng hình học và bề ngoài tương đương hoặc tốt hơn các phương pháp SOTA như NeuS, Instant-NGP, Instant-NSR, Voxurf, D-NeRF, TiNeuVox, đặc biệt là hình học chi tiết và ít nhiễu hơn Instant-NGP \cite{Wang2023NeuS2_source_3, Wang2023NeuS2_source_7, Wang2023NeuS2_source_9, Wang2023NeuS2_source_20, Wang2023NeuS2_source_53, Wang2023NeuS2_source_173, Wang2023NeuS2_source_181, Wang2023NeuS2_source_190, Wang2023NeuS2_source_194, Wang2023NeuS2_source_199, Wang2023NeuS2_source_209}.
    \item \textbf{Khả năng Xử lý Cảnh động Phức tạp:} Kiến trúc huấn luyện tăng cường kết hợp GTP cho phép xử lý hiệu quả các chuỗi video dài (hàng nghìn frame) với chuyển động lớn và biến dạng phức tạp, vượt qua giới hạn của nhiều phương pháp trước đó \cite{Wang2023NeuS2_source_6, Wang2023NeuS2_source_27, Wang2023NeuS2_source_34, Wang2023NeuS2_source_60, Wang2023NeuS2_source_62, Wang2023NeuS2_source_104, Wang2023NeuS2_source_152, Wang2023NeuS2_source_183, Wang2023NeuS2_source_195, Wang2023NeuS2_source_204, Wang2023NeuS2_source_205}.
    \item \textbf{Huấn luyện Ổn định:} Các kỹ thuật như tính đạo hàm bậc hai chính xác và huấn luyện tiến bộ giúp quá trình học ổn định hơn so với các phương pháp xấp xỉ hoặc huấn luyện trực tiếp trên độ phân giải cao \cite{Wang2023NeuS2_source_5, Wang2023NeuS2_source_100, Wang2023NeuS2_source_133, Wang2023NeuS2_source_178}.
\end{itemize}

\subsubsection{Nhược điểm và Hạn chế}
\begin{itemize}
    \item \textbf{Thiếu Đối ứng Dày đặc (Dense Correspondences):} Đối với cảnh động, NeuS2 tái tạo hình học cho từng frame một cách hiệu quả nhưng không cung cấp trực tiếp thông tin đối ứng dày đặc giữa các điểm trên bề mặt qua thời gian \cite{Wang2023NeuS2_source_221}. Việc này có thể cần các bước xử lý hậu kỳ như tracking lưới \cite{Wang2023NeuS2_source_222}.
    \item \textbf{Chi phí Lưu trữ cho Cảnh động:} Do lưu trữ tham số (chủ yếu là bảng băm, khoảng 25MB/frame) cho mỗi frame, việc lưu trữ các chuỗi video rất dài có thể tốn kém \cite{Wang2023NeuS2_source_232}. Cần các kỹ thuật nén trong tương lai \cite{Wang2023NeuS2_source_233}.
    \item \textbf{Phụ thuộc vào Ảnh Đa góc nhìn và Hiệu chỉnh máy ảnh:} Giống như các phương pháp dựa trên NeRF/NeuS khác, NeuS2 yêu cầu đầu vào là ảnh/video từ nhiều góc nhìn với thông số máy ảnh đã biết.
    \item \textbf{Sử dụng Mặt nạ (Masks):} Các thí nghiệm trong bài báo (ví dụ trên DTU) sử dụng mặt nạ đối tượng \cite{Wang2023NeuS2_source_167, Wang2023NeuS2_source_349}. Mặc dù không được nêu rõ là bắt buộc, mặt nạ có thể cần thiết hoặc giúp cải thiện chất lượng/tốc độ hội tụ.
\end{itemize}

Tóm lại, NeuS2 đại diện cho một bước tiến đáng kể trong việc tái tạo bề mặt neural, cân bằng xuất sắc giữa tốc độ, chất lượng và khả năng ứng dụng cho cả cảnh tĩnh và động phức tạp.